\section{\texorpdfstring{Tématické okruhy ke zkoušce}{Tematicke okruhy ke zkousce}}

% =====================================================================================================================
	\begin{frame}
    \frametitle{Číselné soustavy - okruh 10 bodů}
		\small
		\begin{enumerate}
			\item Poziční číselná soustava, celočíselná část, neceločíselná část (\textbf{přesnost!!!}).
			\item Převod do desítkové soustavy, přímé vyjádření a Hornerovo schéma.
			\item Převod do jiné než desítkové soustavy, postupné dělení (násobení) a hledání mocnin.
			\item Převody mezi soustavami, kde je základ jedné mocninou té druhé - proč to lze?
			\item Sčítání (odčítání s přímým vyjádřením znaménka), přetečení, délka operandů a výsledku.
			\item Násobení, příklad pro binární, délka operandů a výsledku.
			\item Přímý kód. Jak se provádějí aritmetické operace?
			\item Aditivní kód. Jak se provádějí aritmetické operace?
			\item Doplňkový kód. Jak se provádějí aritmetické operace?
			\item Pohyblivá řádová čárka IEEE-754, stačí jednoduchá přesnost, ale vědět i speciální hodnoty.
		\end{enumerate}
		
	\end{frame}

% =====================================================================================================================
	\begin{frame}
    \frametitle{Logické obvody - okruh 10 bodů}
		\small
		\begin{enumerate}
			\item Popis logické funkce, tabulka, rovnice, mapa.
			\item Booleova algebra, zákony - vyzkoušíme na úpravě konkrétní rovnice.
			\item Základní logické operace, AND, OR, NOT a XOR, liniové schéma.
			\item De Morganovy zákony, co platí, ukázat na úpravě rovnice.
			\item Booleovská krychle, Karnaughova mapa, tvorba, zápis stavů, minimalizace.
			\item Kombinační obvody, dekodér, multiplexer, demultiplexer, sčítačka.
			\item Sekvenční logická funkce, diagram, logický automat.
			\item Klopné obvody RS, D, zapojení posuvného registru nebo čítače.
			\item Hazardy.
			\item Třístavový výstup a přístup na sběrnici.
			\item ALU, co dělá, příznaky, akumulátor atd.
		\end{enumerate}
		
	\end{frame}
	
% =====================================================================================================================
	\begin{frame}
    \frametitle{Mikroprocesory - okruh 10 bodů}
		\small
		\begin{enumerate}
			\item Paměť RAM, FLASH, EEPROM, reprezentace dat Big a Little Endian.
			\item Typy architektur, von Neumannova a Harvardská, rozdíly, problémy.
			\item Instrukční sada CISC a RISC.
			\item STM8 CPU registry, PC, SP, A, X, Y, CC.
			\item Instrukce a strojový kód, pipeline, instrukce skoků.
			\item Podmínková větvení v assembleru.
			\item Přímá adresace a nepřímá adresace.
			\item Indexová adresace a relativní adresace.
			\item Zásobníková paměť, k čemu slouží, přístup a adresace.
			\item Funkce a přerušení, rozdíly a práce se zásobníkem.
			\item Periferie procesoru, konfigurace, stav po resetu apod.
		\end{enumerate}
		
	\end{frame}

% =====================================================================================================================
	\begin{frame}
    \frametitle{Záchranná otázka - 4 body}
		\small
		\begin{enumerate}
			\item BCD kód.
			\item Parita.
			\item Práce s maskou registru.
			\item Rozdíl mezi asynchronním a synchronním čítačem.
			\item SCADA, k čemu je to dobré.
			\item OPC server, k čemu je to dobré.
			\item Control Web, přístroje, proměnné, procedury, ovladače.
			\item Rozdíl mezi synchronní a asynchronní komunikací.
			\item UART komunikace.
			\item SPI komunikace.
		\end{enumerate}
		
	\end{frame}